\section{Tối ưu hóa quỹ đạo và Điều khiển tối ưu (Optimal Control)}

\subsection{Các phương pháp dựa trên Tối ưu hóa quỹ đạo}
Tối ưu hóa quỹ đạo (Trajectory Optimization) là nhóm phương pháp mạnh mẽ để sinh ra các quỹ đạo mượt mà và khả thi về mặt động lực học. Các phương pháp phổ biến như CHOMP ~\cite{Zucker2013_CHOMP} hay TrajOpt ~\cite{Schulman2014_TrajOpt} mô hình hóa việc tránh va chạm dưới dạng các hàm phạt (penalty functions) và sử dụng quy hoạch phi tuyến (SQP) để tìm nghiệm cục bộ. Mặc dù cho kết quả tốt trong các môi trường đơn giản, các thuật toán này phụ thuộc rất lớn vào giá trị khởi tạo. Nếu điểm khởi tạo nằm trong một cực tiểu cục bộ (ví dụ: bị bao quanh bởi vật cản hình chữ U), thuật toán thường thất bại trong việc tìm ra đường đi hợp lệ.

\subsection{Áp dụng Điều khiển tối ưu (OCP) và Multiple Shooting}
Để khắc phục tính phi lồi tự nhiên của không gian có vật cản, bài toán hoạch định chuyển động thường được phát biểu dưới dạng Bài toán Điều khiển Tối ưu (Optimal Control Problem). Tuy nhiên, để giải quyết chúng trong thời gian thực, các kỹ thuật rời rạc hóa phải được áp dụng một cách cẩn thận. 

Phương pháp Single Shooting (mô phỏng tiến từ trạng thái ban đầu) rất đơn giản nhưng cực kỳ nhạy cảm với sai số số học và dễ dẫn đến sự bất ổn định phi tuyến, đặc biệt với các hệ thống robot có thời gian di chuyển dài ~\cite{BockPlitt1984}. Ngược lại, phương pháp Multiple Shooting ~\cite{DiehlEtAl2006} đã chứng minh được tính ưu việt bằng cách chia nhỏ quỹ đạo thành các đoạn ngắn và giải quyết đồng thời. Khả năng tích hợp Multiple Shooting với các bộ giải quy hoạch phi tuyến hiệu năng cao (như IPOPT hay SNOPT) đã giúp phương pháp này trở thành tiêu chuẩn vàng trong điều khiển tối ưu. Mặc dù vậy, khi áp dụng thuần túy vào môi trường có nhiều vật cản, Multiple Shooting vẫn không thể vượt qua hoàn toàn rào cản về cực tiểu cục bộ nếu không có sự phân chia không gian an toàn một cách có hệ thống.

Sự kết hợp giữa lý thuyết đồ thị (như GCS) để tìm một đường hầm (tube) không gian an toàn lồi, sau đó áp dụng Multiple Shooting trong không gian đó, đã được đề xuất trong các công trình gần đây nhằm tận dụng tối đa ưu điểm của cả hai: sự đảm bảo về an toàn hình học toàn cục và sự chính xác, khả thi về mặt động lực học của hệ thống liên tục.
